%%%%%%%%%%%%%%%%%%%%%%%%%%%%%%%%%%%%%%%%%%%%%%%%%%
\needspace{12\baselineskip}
\section{Java: Arrays und Aufzählungen}
%%%%%%%%%%%%%%%%%%%%%%%%%%%%%%%%%%%%%%%%%%%%%%%%%%
\worum
Bisher enthielt jede Variable genau einen Wert. Ein \textbf{Array} (Feld) fasst viele gleichartige Werte unter einem Namen zusammen; auf die einzelnen Elemente greift man über einen \textbf{Index} zu. Arrays sind in Java \textbf{Objekte}: Die Array-Variable liegt auf dem Stack und enthält nur eine \textbf{Referenz} (Adresse) auf das eigentliche Array, das mit \code{new} im Heap angelegt wird. Das hat eine wichtige Folge: Weist man eine Array-Variable einer anderen zu, wird nicht das Array kopiert, sondern nur die Referenz.

Ein \textbf{Aufzählungstyp} (\code{enum}) ist ein Datentyp, der als Werte nur eine feste Menge benannter Konstanten zulässt.

\begin{merke}
\begin{itemize}[itemsep=0.1em]
\item Anlegen: \code{int[] f = new int[5];} \ oder mit Werten: \ \code{int[] f = \{10, 4, 7, -1, 8\};}
\item Der Index läuft von \code{0} bis \code{f.length - 1}. Ein Zugriff außerhalb löst zur Laufzeit eine \code{ArrayIndexOutOfBoundsException} aus.
\item Die Länge steht in \code{f.length} und lässt sich nach dem Anlegen nicht mehr ändern.
\item Mit \code{new} angelegte Arrays sind vorbelegt: Zahlen mit \code{0}, \code{boolean} mit \code{false}, Referenzen (z.\,B. \code{String}) mit \code{null}.
\item Durchlaufen: mit Index \code{for (int i = 0; i < f.length; i++)} (Lesen und Ändern) oder mit der For-each-Schleife \code{for (int wert : f)} (nur Lesen).
\item \textbf{Mehrdimensional:} \code{int[][] m = new int[2][3];} ist ein Array aus 2 Arrays mit je 3 Elementen. \code{m.length} ist 2, \code{m[0].length} ist 3. Die Zeilen dürfen verschieden lang sein.
\item \textbf{Enum:} \code{enum AmpelFarbe \{ ROT, GELB, GRUEN \}}, Verwendung: \code{AmpelFarbe a = AmpelFarbe.ROT;}
\end{itemize}
\end{merke}

\beispiele
\beispiel{Summe und Maximum}
\begin{javacode}[Ausschnitt]{Durchlaufen}
int[] werte = {10, 4, 7, -1, 8};

int summe = 0;
for (int w : werte) {               // For-each: nur lesen
	summe += w;
}

int max = werte[0];
for (int i = 1; i < werte.length; i++) {
	if (werte[i] > max) {
		max = werte[i];
	}
}
System.out.println(summe + " " + max);      // 28 10
\end{javacode}

\beispiel{Referenzen: zwei Namen für dasselbe Array}
\noindent\begin{minipage}[c]{0.5\linewidth}
\begin{javacode}[Ausschnitt]{Referenz}
int[] a = {1, 2, 3};
int[] b = a;          // kopiert nur die Referenz
b[0] = 99;
System.out.println(a[0]);     // 99
\end{javacode}
\end{minipage}\hfill
\begin{minipage}[c]{0.45\linewidth}
\centering
\begin{tikzpicture}[syntax, font=\small\ttfamily]
	\node[draw=diagramm, semithick, minimum width=9mm, minimum height=6mm, label=left:a] (a) at (0,0.5) {};
	\node[draw=diagramm, semithick, minimum width=9mm, minimum height=6mm, label=left:b] (b) at (0,-0.5) {};
	\foreach \x/\v in {0/99, 1/2, 2/3}
		\node[flaeche, minimum width=9mm, minimum height=6mm] (f\x) at (2.6+\x*0.9,0) {\v};
	\fill[diagramm] (a.center) circle (1.5pt);
	\fill[diagramm] (b.center) circle (1.5pt);
	\draw[linie, ->] (a.center) -- (f0.west);
	\draw[linie, ->] (b.center) -- (f0.west);
	\node[font=\footnotesize\rmfamily] at (0,-1.2) {Stack};
	\node[font=\footnotesize\rmfamily] at (3.5,-1.2) {Heap};
\end{tikzpicture}
\end{minipage}

\smallskip
Eine echte Kopie braucht ein neues Array, in das die Werte einzeln übertragen werden:
\begin{javacode}[Ausschnitt]{Kopie}
int[] c = new int[a.length];
for (int i = 0; i < a.length; i++) {
	c[i] = a[i];
}
\end{javacode}
Aus demselben Grund vergleicht \code{a == b} nicht die Inhalte, sondern nur, ob beide auf \emph{dasselbe} Array zeigen.

\beispiel{Zweidimensionales Array}
\begin{javacode}[Ausschnitt]{Matrix}
int[][] m = { {1, 2, 3}, {4, 5, 6} };       // 2 Zeilen, 3 Spalten
System.out.println(m[1][2]);                // 6 (Zeile 1, Spalte 2)

for (int i = 0; i < m.length; i++) {        // alle Zeilen
	int zeilensumme = 0;
	for (int j = 0; j < m[i].length; j++) { // alle Spalten der Zeile i
		zeilensumme += m[i][j];
	}
	System.out.println("Zeile " + i + ": " + zeilensumme);      // 6 und 15
}
\end{javacode}

\beispiel{Zählen mit einem Hilfsarray}
Ein Würfel wird 600-mal geworfen. Ein Array merkt sich, wie oft jede Augenzahl gefallen ist. Die Augenzahl dient dabei direkt als Index.
\begin{javacode}[Ausschnitt]{Würfelstatistik}
int[] anzahl = new int[7];                  // Index 1 bis 6 wird benutzt, Index 0 bleibt leer
for (int i = 0; i < 600; i++) {
	int wurf = (int) (Math.random() * 6) + 1;   // Zufallszahl von 1 bis 6
	anzahl[wurf]++;
}
for (int augen = 1; augen <= 6; augen++) {
	System.out.println(augen + ": " + anzahl[augen]);
}
\end{javacode}
\code{Math.random()} liefert einen zufälligen \code{double}-Wert $x$ mit $0 \leq x < 1$.

\beispiel{Aufzählungstyp mit switch}
\begin{javacode}{EnumTest.java}
public class EnumTest {
	enum AmpelFarbe { ROT, GELB, GRUEN }

	public static void main(String[] args) {
		AmpelFarbe ampel = AmpelFarbe.GELB;
		switch (ampel) {
			case ROT:   System.out.println("Stehen bleiben"); break;
			case GELB:  System.out.println("Achtung"); break;
			case GRUEN: System.out.println("Gehen"); break;
		}
	}
}
\end{javacode}

\aufgaben
\aufgabe{Arrays füllen}
\begin{enumerate}[label=\alph*)]
\item Deklarieren Sie ein \code{int}-Array mit 200 Elementen und füllen Sie es mit den ersten 200 ungeraden Zahlen (1, 3, 5, \ldots).
\item Deklarieren Sie ein \code{int}-Array mit 50 Elementen und füllen Sie es mit den ersten 50 Quadratzahlen in absteigender Reihenfolge (2500, 2401, \ldots, 1).
\item Füllen Sie ein Array mit den ersten 30 Fibonacci-Zahlen (0, 1, 1, 2, 3, 5, 8, \ldots). Jede Zahl ist die Summe ihrer beiden Vorgänger.
\end{enumerate}

\aufgabe{Lässt sich das ausführen?}
Lassen sich die folgenden Ausschnitte ausführen? Falls nein: Tritt der Fehler beim Übersetzen oder zur Laufzeit auf? Falls ja: Was wird ausgegeben?
\begin{tcbraster}[raster columns=2, raster equal height=rows, raster column skip=3mm, raster row skip=1mm]
\begin{javacode}[Ausschnitt]{a)}
int[] f = new int[3];
System.out.println(f[1]);
\end{javacode}
\begin{javacode}[Ausschnitt]{b)}
int[] f = {1, 2, 3};
System.out.println(f[3]);
\end{javacode}
\begin{javacode}[Ausschnitt]{c)}
int[] f;
f[0] = 1;
System.out.println(f[0]);
\end{javacode}
\begin{javacode}[Ausschnitt]{d)}
String[] s = new String[2];
System.out.println(s[0]);
\end{javacode}
\begin{javacode}[Ausschnitt]{e)}
int[] f = {4, 5, 6};
f[f.length - 1] = f[0];
System.out.println(f[2] + f.length);
\end{javacode}
\begin{javacode}[Ausschnitt]{f)}
int[][] m = new int[2][];
m[0] = new int[3];
m[1] = new int[4];
System.out.println(m.length + m[1].length);
\end{javacode}
\end{tcbraster}

\aufgabe[\K]{Programmverständnis}
Wie lautet die Ausgabe des folgenden Programmausschnitts? Erläutern Sie Ihre Antwort kurz, am besten mit einer Skizze von Stack und Heap.
\begin{javacode}[Ausschnitt]{Referenzen}
int[] a = { 5, 6 };
int[] b = { 7, 8, 9 };
int[] c = a;

a = b;
a[2] = 1;
c[0] += b[2];

for (int i = 0; i < c.length; i++) {
	System.out.println("a[" + i + "]=" + a[i] + " b[" + i + "]=" + b[i] + " c[" + i + "]=" + c[i]);
}
\end{javacode}

\aufgabe{Programme mit Arrays}
\begin{enumerate}[label=\alph*)]
\item Bestimmen Sie in einem \code{double}-Array das Minimum und die Position, an der es steht.
\item Berechnen Sie den Mittelwert der Elemente und zählen Sie, wie viele Elemente über dem Mittelwert liegen.
\item Drehen Sie die Reihenfolge der Elemente eines Arrays um, ohne ein zweites Array zu verwenden.
\item Prüfen Sie, ob ein Array aufsteigend sortiert ist.
\end{enumerate}

\aufgabe{Zweidimensionale Arrays}
\begin{enumerate}[label=\alph*)]
\item Legen Sie ein Array \code{int[][] einmaleins = new int[10][10]} an und füllen Sie es so, dass \code{einmaleins[i][j]} das Produkt $(i+1) \cdot (j+1)$ enthält. Geben Sie es zeilenweise aus.
\item Berechnen Sie für eine quadratische Matrix die Summe der Hauptdiagonalen.
\item Ein Notenspiegel speichert für 4 Studierende je 3 Noten in einem \code{double[4][3]}. Geben Sie für jede Person den Durchschnitt aus und anschließend den Durchschnitt jeder der 3 Prüfungen.
\end{enumerate}

\aufgabe{Aufzählungen}
Legen Sie einen Aufzählungstyp \code{Wochentag} mit den Werten \code{MO} bis \code{SO} an. Schreiben Sie ein Programm, das für eine Variable vom Typ \code{Wochentag} mit einer \code{switch}-Anweisung ausgibt, ob es sich um einen Werktag oder um das Wochenende handelt. Welchen Vorteil hat der Aufzählungstyp gegenüber einer Codierung der Tage durch die Zahlen 1 bis 7?

\aufgabe[\S]{Pascalsches Dreieck}
Im pascalschen Dreieck stehen in der $i$-ten Zeile genau $i$ Werte. Am Rand steht jeweils eine 1, jeder andere Wert ist die Summe der beiden darüber liegenden Werte. Legen Sie ein \emph{ungleichmäßiges} Array \code{int[][] pascal = new int[n][]} an, berechnen Sie die ersten $n$ Zeilen und geben Sie sie aus.
\begin{konsole}
1
1 1
1 2 1
1 3 3 1
1 4 6 4 1
\end{konsole}

\ueben
\aufgabe{Arrays anlegen}
Schreiben Sie jeweils die Anweisung zum Anlegen des Arrays.
\begin{enumerate}[label=\alph*)]
\item ein Array mit den Namen der zwölf Monate
\item ein Array für 100 Messwerte (Kommazahlen)
\item ein Spielfeld aus $3 \times 3$ Zeichen
\item ein Array, das für die Zahlen 0 bis 99 speichert, ob sie Primzahlen sind
\end{enumerate}

\aufgabe{Was wird ausgegeben?}
\begin{tcbraster}[raster columns=2, raster equal height=rows, raster column skip=3mm, raster row skip=1mm]
\begin{javacode}[Ausschnitt]{a)}
int[] f = {3, 1, 4, 1, 5};
int s = 0;
for (int i = 0; i < f.length; i += 2) {
	s += f[i];
}
System.out.println(s);
\end{javacode}
\begin{javacode}[Ausschnitt]{b)}
int[] f = {3, 1, 4, 1, 5};
for (int i = 1; i < f.length; i++) {
	f[i] = f[i] + f[i - 1];
}
System.out.println(f[4]);
\end{javacode}
\begin{javacode}[Ausschnitt]{c)}
int[] f = {1, 2, 3};
for (int w : f) {
	w = w * 2;
}
System.out.println(f[0] + f[1] + f[2]);
\end{javacode}
\begin{javacode}[Ausschnitt]{d)}
int[] a = {1, 2};
int[] b = {1, 2};
int[] c = a;
System.out.println(a == b);
System.out.println(a == c);
\end{javacode}
\end{tcbraster}

\aufgabe{Programme mit eindimensionalen Arrays}
\begin{enumerate}[label=\alph*)]
\item Zählen Sie, wie viele Elemente eines \code{int}-Arrays positiv sind.
\item Bestimmen Sie das zweitgrößte Element.
\item Verschieben Sie alle Elemente um eine Position nach links; das erste Element wandert ans Ende.
\item Addieren Sie zwei gleich lange Arrays elementweise in ein drittes Array.
\item Prüfen Sie, ob zwei Arrays denselben Inhalt haben.
\item Erzeugen Sie eine Kopie eines Arrays in umgekehrter Reihenfolge.
\end{enumerate}

\aufgabe{Programme mit zweidimensionalen Arrays}
\begin{enumerate}[label=\alph*)]
\item Erzeugen Sie eine $n \times n$-Matrix, die auf der Hauptdiagonalen Einsen und sonst Nullen enthält.
\item Bestimmen Sie für eine Matrix das Maximum jeder Zeile und das Maximum jeder Spalte.
\item Erzeugen Sie zu einer Matrix mit $m$ Zeilen und $n$ Spalten die \emph{transponierte} Matrix mit $n$ Zeilen und $m$ Spalten (Zeilen und Spalten sind vertauscht).
\item Ein Tic-Tac-Toe-Feld ist als \code{char[3][3]} gespeichert. Prüfen Sie, ob in einer Zeile dreimal dasselbe Zeichen (\code{'X'} oder \code{'O'}) steht.
\end{enumerate}

\aufgabe{Himmelsrichtungen}
Legen Sie den Aufzählungstyp \code{Richtung} mit den Werten \code{NORD}, \code{OST}, \code{SUED} und \code{WEST} an. Schreiben Sie ein Programm, das zu einer Richtung mit \code{switch} die Richtung bestimmt, die sich nach einer Vierteldrehung nach rechts ergibt.
